\documentclass{article}
\usepackage[final]{neurips_2024}

\usepackage[utf8]{inputenc}
\usepackage[T1]{fontenc}
\usepackage{hyperref}
\usepackage{url}
\usepackage{booktabs}
\usepackage{amsfonts}
\usepackage{nicefrac}
\usepackage{microtype}
\usepackage{xcolor}
\usepackage{graphicx}
\usepackage{amsmath}
\usepackage{multirow}

\title{Degree-Aware Curriculum Hard-Negative Sampling for LightGCL}

\author{
  Ngo Hoang Dang \quad Mai Duc Phuc \quad Nguyen Hoang Dang \quad Nguyen Kim Tuyen \quad Mai Hoang Long\\
  Graph Analytics for Big Data (IT5429E) --- Class Project Report
}

\begin{document}

\maketitle

\begin{abstract}
LightGCL (ICLR 2023) is a graph contrastive learning model for recommendation that
contrasts a LightGCN view of the user--item graph with an SVD-reconstructed view, yet its
ranking objective still relies on \emph{uniformly} sampled negative items in the BPR loss.
We propose a degree-aware curriculum hard-negative sampler for this BPR component:
negatives are drawn from a mixture of a uniform distribution and a Personalized-PageRank
(PPR) hard-negative distribution, whose mixture weight grows with training progress
(curriculum) and with each user's activity level (degree), so that sparse users keep
mostly random negatives, limiting false-negative bias. On Yelp2018, a fixed-ramp mixture
(S2) and our degree-gated curriculum mixture (S3) both improve over the uniform baseline
(S0) and over pure PPR hard negatives (S1) on Recall@20/NDCG@20 and on every degree
bucket (head/mid/tail), consistently across 3 seeds; however, S3's extra degree-gating
does not show a measurable additional benefit over the simpler S2 in our experiments,
including in the tail bucket where it was expected to matter most. SimGCL, an external
contrastive-learning baseline, outperforms all LightGCL variants tested. We report all
numbers honestly, including this partly-negative finding on the core research question.
\end{abstract}

\section{Introduction}

Graph-based collaborative filtering models such as LightGCN~\cite{lightgcn} and its
contrastive-learning successor LightGCL~\cite{lightgcl} rank items for a user via an
inner product of learned embeddings, trained with the pairwise BPR loss~\cite{bpr}. BPR
requires, for every observed (user, item) interaction, a sampled \emph{negative} item
that the user is assumed not to prefer. In nearly all implementations, including
LightGCL's official code, this negative is drawn \emph{uniformly at random} from the
user's non-interacted items. On sparse implicit-feedback graphs this is a reasonable
default but a known missed opportunity: most uniformly-sampled negatives are "easy"
(unrelated to the user's taste) and contribute little gradient signal, while a small
number of \emph{hard} negatives --- items structurally close to the user in the
interaction graph but never interacted with --- would push the model harder. The risk of
hard-negative mining is \emph{false-negative bias}: an item that looks hard precisely
because it is popular among the user's neighborhood is disproportionately likely to be an
unobserved \emph{positive}, especially for users with very few training interactions,
where the model has the least evidence to tell the two apart.

\textbf{Research question.} Does a hard-negative curriculum modulated by user activity
improve the recommendation accuracy and long-tail robustness of LightGCL without
increasing false-negative bias?

We answer this with a controlled ablation, changing \emph{only} the BPR negative sampler
of the official LightGCL implementation and leaving its SVD contrastive view and
InfoNCE loss untouched, then compare four sampler variants (uniform, pure hard-negative,
fixed-ramp mixture, and degree-gated curriculum mixture) against each other and against
three external baselines (MF-BPR, LightGCN, SimGCL) on Yelp2018.

\section{Application domain: dataset, task, and metrics}

\paragraph{Dataset.} We use Yelp2018, an implicit-feedback recommendation benchmark of
user check-ins at local businesses, using the exact train/test splits shipped with the
official LightGCL repository~\cite{lightgcl}. \textbf{Note:} the statistics below are
measured directly from that shipped data, which differs from our own project proposal
draft's Table~1 (31{,}668 users / 38{,}048 items) --- see Appendix~\ref{app:dataset-note}
for the discrepancy and why we use the measured numbers.

\begin{table}[h]
\centering
\caption{Statistics of the datasets actually used (measured from the shipped splits).}
\label{tab:dataset-stats}
\begin{tabular}{lrrrrr}
\toprule
Dataset & \#Users & \#Items & Train & Test & Density \\
\midrule
Yelp2018 (primary) & 29{,}601 & 24{,}734 & 1{,}069{,}128 & 305{,}466 & 0.188\% \\
Gowalla (PPR precomputed; not evaluated, see \S\ref{sec:setup}) & 50{,}821 & 57{,}440 & 1{,}172{,}425 & 130{,}270 & 0.045\% \\
\bottomrule
\end{tabular}
\end{table}

\paragraph{Task and metrics.} Top-$K$ recommendation from implicit feedback: given the
bipartite user--item interaction graph, rank all items a user has not interacted with and
recommend the top $K$. We report Recall@20 and NDCG@20 under the standard full-ranking
protocol (rank all non-train items, following the official LightGCL evaluation code),
averaged over test users, plus a breakdown by user-degree bucket (head/mid/tail, tertiles
by training-interaction count; thresholds: tail~$\le 20$, mid~$\le 33$, head~$>33$
interactions).

\section{Method}
\label{sec:method}

\subsection{Base model and BPR objective}
Users $u$ and items $i$ are embedded as $e_u, e_i \in \mathbb{R}^d$, scored by inner
product $\hat y_{ui} = e_u^\top e_i$. LightGCN~\cite{lightgcn} refines embeddings with
$L$ layers of symmetric-normalized linear message passing over the bipartite graph and
averages all layers' representations. Training uses the pairwise BPR loss: for each
observed pair $(u,i)$, a negative item $j$ is sampled and the model is pushed to rank $i$
above $j$,
\[
\mathcal{L}_{\mathrm{bpr}} = \sum_{(u,i)\in O} \mathbb{E}_{j\sim P(\cdot\mid u)}\big[-\log
\sigma(\hat y_{ui} - \hat y_{uj})\big],
\]
where $P(\cdot\mid u)$ is uniform over non-interacted items in all baselines. This
sampler is the only component we modify.

\subsection{LightGCL: SVD-guided graph contrastive learning}
LightGCL~\cite{lightgcl} augments LightGCN with a second view built from a truncated
rank-$q$ SVD of the normalized adjacency, and aligns the two views with an InfoNCE loss
using in-batch negatives, added to $\mathcal{L}_{\mathrm{bpr}}$ with weight
$\lambda_1$. Critically, we verified in the official code (\texttt{model.py}) that the
\emph{ranking score itself} is computed purely from the LightGCN branch; the SVD view
only ever feeds the contrastive loss term. This means LightGCN and MF-BPR baselines can
be obtained for free as ablations of the same codebase (\S\ref{sec:setup}), and that our
sampler change --- which only affects $\mathcal{L}_{\mathrm{bpr}}$'s negative $j$ ---
cannot alter the contrastive branch at all.

\subsection{Proposed: degree-aware curriculum hard-negative sampling}
Following PinSAGE~\cite{pinsage}, we rank all non-interacted items for each user by
Personalized PageRank (PPR) on the interaction graph and treat items in a percentile band
of that ranking as hard-negative candidates,
\[
B_u = \{\, j \notin N_u : \mathrm{rank}_{\mathrm{PPR}}(j\mid u) \in [p_{lo}|I|,\,
p_{hi}|I|] \,\}, \qquad P_{\mathrm{PPR}}(j\mid u) = \mathrm{Unif}(B_u).
\]
At training epoch $t$, the BPR negative for user $u$ is drawn from a mixture
\[
P_t(j\mid u) = (1-\alpha_{u,t})\, P_{\mathrm{uniform}}(j) + \alpha_{u,t}\, P_{\mathrm{PPR}}(j\mid u),
\]
where the mixture weight combines a curriculum ramp with a degree gate,
\[
\alpha_{u,t} = s(t)\cdot \sigma\big(a\,(\log(1+d_u) - \log(1+d_{\mathrm{mid}}))\big), \qquad s(t) = \min(1, t/T_w).
\]
$s(t)$ ramps hard negatives on only after a $T_w$-epoch warm-up (curriculum learning,
\cite{curriculum}); the sigmoid gate gives high-degree (active) users harder negatives
earlier while low-degree (sparse) users keep mostly uniform negatives throughout, since a
PPR-similar item is more likely to be an unobserved positive precisely where a sparse
user gives the model the least evidence. We compare four sampler modes, all applied only
to LightGCL's BPR term:
\begin{itemize}
\item \textbf{S0} --- uniform (the unmodified baseline).
\item \textbf{S1} --- pure PPR hard negatives ($\alpha_{u,t}\equiv 1$, no curriculum/gate).
\item \textbf{S2} --- fixed global mixture ($\alpha_{u,t} = s(t)\cdot\bar\alpha$, no degree
dependence).
\item \textbf{S3} --- our full method: degree-gated curriculum mixture, as above.
\end{itemize}
Full implementation details (GPU-batched PPR via power iteration, percentile-band
construction, vectorized per-epoch sampling) are in
\texttt{docs/design/sampler\_design.md} and \texttt{scripts/precompute\_ppr.py}; PPR is
precomputed once per dataset (11.4s for Yelp2018, 30.1s for Gowalla on an RTX 5090 ---
both far under the design budget) and cached as a fixed-size candidate pool per user.

\section{Experimental setup}
\label{sec:setup}

\paragraph{Baselines.} MF-BPR, LightGCN, SimGCL, and LightGCL. MF-BPR and LightGCN are
obtained as ablations of the same LightGCL codebase (\texttt{--gnn\_layer 0 --lambda1 0}
and \texttt{--lambda1 0} respectively; see \S\ref{sec:method}). SimGCL~\cite{simgcl} is
implemented separately (\texttt{scripts/baselines/simgcl.py}) since it requires
noise-perturbation-based contrastive augmentation rather than an SVD view; it uses its
own (unmodified, uniform) negative sampler, matching its official design.

\paragraph{Ablation.} S0/S1/S2/S3 as defined in \S\ref{sec:method}, all on LightGCL, on
Yelp2018.

\paragraph{Protocol.} All models: $d=64$, 2 GNN layers, 100 epochs, Adam, following
LightGCL's default architecture hyperparameters unless noted below. \emph{Screening}
(single seed) swept the PPR band $(p_{lo}, p_{hi}) \in \{(0.05,0.15), (0.02,0.10),
(0.10,0.20)\}$ for S1, selecting the best band $(0.02,0.10)$ for S2/S3; $\bar\alpha \in
\{0.3, 0.5\}$ for S2; and $(a, T_w) \in \{(1.0,10),(1.5,10),(1.0,5)\}$ for S3 --- 9
LightGCL-ablation runs plus MF-BPR/LightGCN screening, 11 runs total, all on a single
RTX 5090, none diverged (see \S\ref{sec:stability}). \emph{Final} numbers re-run the best
band/hyperparameters for S0, S2 ($\bar\alpha=0.3$), and S3 ($a=1.0, T_w=10$) with two
additional seeds (3 seeds total per config, with per-user metrics saved for the
degree-bucket table), and MF-BPR/LightGCN/S1-best with one additional seed (2 seeds
total); SimGCL was run fresh at 2 seeds. Time budget did not permit a Gowalla run or the
secondary SVD-rank ablation from the proposal; both are noted as future work.

\paragraph{Implementation notes / stability fixes.}
\label{sec:stability}
Two numerical-stability bugs in the unmodified LightGCL reference code caused training to
diverge to \texttt{NaN} on long (100-epoch) runs, discovered while reproducing the
MF-BPR and LightGCN ablations:
\begin{enumerate}
\item The BPR loss computed \texttt{sigmoid(x).log()} directly, which returns $-\infty$
once score margins grow large enough to saturate the sigmoid in float32. We replaced it
with the numerically-stable \texttt{F.logsigmoid(x)} and added a defensive
\texttt{clamp(x, -20, 20)}, which is a no-op for well-behaved training but blocks
runaway growth (the clamp's zero gradient outside the bounds also acts as an implicit
stabilizer).
\item The InfoNCE-style contrastive loss computed \texttt{log(exp(x).sum())} without the
log-sum-exp trick; once embeddings grow, \texttt{exp(x)} overflows to \texttt{inf}, and
\texttt{log(inf)=inf} poisons the total loss \emph{even when $\lambda_1=0$} (i.e. even
for the LightGCN ablation, which does not use this term for anything but still computes
it), since $0 \times \mathrm{NaN} = \mathrm{NaN}$. We replaced it with
\texttt{torch.logsumexp}, mathematically identical for finite inputs but overflow-safe.
\end{enumerate}
We also added global gradient-norm clipping (max\_norm=5) and, since the reference code
has no seed control and its \texttt{--decay} flag is parsed but never applied to the
optimizer, added explicit \texttt{--seed} plumbing (torch + numpy) for reproducibility.
After these fixes, we stress-tested LightGCN for a full 100 epochs and observed
\emph{zero} \texttt{NaN} occurrences (final Recall@20=0.0898); before the fixes, the same
configuration reproducibly diverged between epochs 9 and 93 depending on random
initialization. MF-BPR additionally needed a lower learning rate and higher $L_2$ weight
($\mathrm{lr}=10^{-4}$, $\lambda_2=10^{-5}$ vs. the repository's LightGCN-tuned
$\mathrm{lr}=10^{-3}$, $\lambda_2=10^{-7}$) since it has no graph-smoothing to bound
embedding scale. These fixes are orthogonal to our sampler contribution but were
necessary for \emph{any} of the baselines to produce a trustworthy 100-epoch number, and
we release them as part of this codebase.

\section{Results}

\begin{table}[h]
\centering
\caption{Main comparison on Yelp2018, Recall@20/NDCG@20, mean$\pm$std across seeds.}
\label{tab:main-results}
\begin{tabular}{lrcc}
\toprule
Model & \#seeds & Recall@20 & NDCG@20 \\
\midrule
MF-BPR & 2 & 0.0299 $\pm$ 0.0002 & 0.0268 $\pm$ 0.0001 \\
LightGCN & 2 & 0.0901 $\pm$ 0.0000 & 0.0764 $\pm$ 0.0001 \\
LightGCL (S0, uniform baseline) & 3 & 0.1005 $\pm$ 0.0008 & 0.0866 $\pm$ 0.0007 \\
LightGCL + S3 (ours, degree-gated curriculum) & 3 & 0.1015 $\pm$ 0.0006 & 0.0875 $\pm$ 0.0003 \\
LightGCL + S2 (fixed-ramp mixture) & 3 & \textbf{0.1019} $\pm$ 0.0007 & \textbf{0.0879} $\pm$ 0.0005 \\
SimGCL & 2 & \textbf{0.1080} $\pm$ 0.0002 & \textbf{0.0933} $\pm$ 0.0001 \\
\bottomrule
\end{tabular}
\end{table}

Table~\ref{tab:main-results} shows the expected ordering among the LightGCL-family
models: MF-BPR $<$ LightGCN $<$ LightGCL(S0), confirming graph propagation and then
contrastive learning each add value on this dataset, and our reproduction is internally
consistent. Both hard-negative mixture variants (S2, S3) improve over the uniform
baseline S0, with S3 (our proposed method) beating S0 on \emph{every one of the 3
seeds individually} (0.1009/0.1021/0.1016 vs. 0.1002/0.1015/0.0999) --- a small but
consistent effect. Unexpectedly, SimGCL --- an external baseline whose own negative
sampler we did not touch --- outperforms every LightGCL variant we tested by a wide
margin. We did not have time to verify whether this gap is a genuine model-quality
difference or an artifact of hyperparameter tuning (SimGCL's noise scale $\epsilon$ and
CL weight were left at commonly-used defaults, not tuned for this dataset); we report it
as an honest, unresolved observation rather than adjust after the fact.

\subsection{Ablation: sampling strategy}

\begin{table}[h]
\centering
\caption{BPR sampler ablation on LightGCL, Yelp2018 (best band/hyperparameters per
family; screening seed=0 only for single-seed rows).}
\label{tab:ablation}
\begin{tabular}{llrcc}
\toprule
ID & Sampler & \#seeds & Recall@20 & NDCG@20 \\
\midrule
S0 & uniform (baseline) & 3 & 0.1005 $\pm$ 0.0008 & 0.0866 $\pm$ 0.0007 \\
S1 & pure PPR, band $(0.02,0.10)$ & 2 & 0.1000 $\pm$ 0.0009 & 0.0858 $\pm$ 0.0007 \\
S2 & mixture, $\bar\alpha{=}0.3$ & 3 & \textbf{0.1019} $\pm$ 0.0007 & \textbf{0.0879} $\pm$ 0.0005 \\
S3 & degree-gated curriculum (ours) & 3 & 0.1015 $\pm$ 0.0006 & 0.0875 $\pm$ 0.0003 \\
\bottomrule
\end{tabular}
\end{table}

Table~\ref{tab:ablation} shows the pattern the design anticipated in one respect and not
in another. \textbf{As hypothesized:} S1 (pure PPR hard negatives, no curriculum, no
degree protection) is the \emph{only} variant that underperforms the uniform baseline
S0, consistent with the false-negative-bias mechanism the proposal targets --- treating
PPR-similar items as negatives from epoch 0 for every user, including sparse ones with
little evidence, does measurably hurt. \textbf{Not hypothesized:} once a curriculum ramp
is added, the simpler global mixture S2 matches or slightly exceeds our full degree-gated
method S3 on every metric. We discuss this in \S\ref{sec:discussion}.

\subsection{Degree-bucket breakdown}

\begin{table}[h]
\centering
\caption{Recall@20 / NDCG@20 by user-degree bucket (tertiles: tail $\le$20, mid $\le$33,
head $>$33 training interactions), mean over 2 seeds (seeds 1--2, the only runs with
per-user logging).}
\label{tab:bucket}
\begin{tabular}{llcc}
\toprule
Bucket & Model & Recall@20 & NDCG@20 \\
\midrule
\multirow{3}{*}{Tail (sparse)}
 & S0 & 0.0989 & 0.0682 \\
 & S3 (ours) & 0.1011 & 0.0691 \\
 & S2 & \textbf{0.1017} & \textbf{0.0695} \\
\midrule
\multirow{3}{*}{Mid}
 & S0 & 0.1083 & 0.0830 \\
 & S3 (ours) & 0.1089 & 0.0836 \\
 & S2 & \textbf{0.1095} & \textbf{0.0842} \\
\midrule
\multirow{3}{*}{Head (active)}
 & S0 & 0.0944 & 0.1098 \\
 & S3 (ours) & 0.0949 & 0.1110 \\
 & S2 & \textbf{0.0952} & \textbf{0.1113} \\
\bottomrule
\end{tabular}
\end{table}
(Table generated by \texttt{scripts/analyze\_degree\_buckets.py} from
\texttt{results/degree\_bucket\_summary.csv}.)

This is the core long-tail-robustness result, and it is a genuine mixed finding. Both S2
and S3 improve over S0 in \emph{every} bucket, including the tail --- so the central
claim that a hard-negative curriculum can help sparse users \emph{without} net harm is
supported: nowhere does adding hard negatives make the tail bucket worse than uniform
sampling. However, the degree gate's specific mechanism --- suppressing hard negatives
more aggressively for low-degree users --- was designed to give S3 a \emph{larger}
relative gain in the tail bucket than S2. We do not observe that: S2's tail-bucket gain
over S0 ($+0.0028$ Recall@20) is if anything slightly larger than S3's ($+0.0022$).

\section{Discussion}
\label{sec:discussion}

Returning to the research question --- does a hard-negative curriculum modulated by user
activity improve accuracy and long-tail robustness without increasing false-negative
bias? --- our results support the first half more clearly than the second. A curriculum
mixture of uniform and PPR hard negatives (S2 or S3) improves both overall accuracy and
every degree bucket over pure uniform sampling (S0), and pure hard-negative sampling
without any curriculum or protection (S1) is the one variant that backfires, exactly as
the false-negative-bias hypothesis predicts. But the \emph{degree-specific} gating that
is this project's distinguishing contribution over a simpler fixed-ramp mixture does not
show a measurable edge in our experiments, including specifically in the tail bucket
where we expected it to matter most.

We see three plausible, non-exclusive explanations, none of which we had time to test
directly: (1) \emph{effect size vs. seed noise} --- differences between S2 and S3 are
$\le 0.0006$ Recall@20 with per-config std around $0.0006$--$0.0008$ from only 2--3
seeds, so the two methods may simply be statistically indistinguishable at this sample
size, and a larger seed count could reverse or confirm the ranking; (2) the degree gate's
hyperparameters ($a$, $d_{\mathrm{mid}}$) were only lightly swept (2 values of $a$, one
$T_w$ reduction) in a single-seed screening pass, so the gate's shape may not be
well-tuned; (3) Yelp2018 at $0.188\%$ density with a median degree around 20 may not have
enough separation between "sparse" and "active" users for the degree gate to meaningfully
differ from a flat mixture --- the proposal's planned second dataset, Gowalla (sparser,
$0.045\%$ density), might show a larger gap, but we did not have time to run it.

We also did not anticipate SimGCL's large margin over every LightGCL variant
(Table~\ref{tab:main-results}); since we did not tune SimGCL's hyperparameters for this
dataset, we cannot yet attribute this to model quality vs. tuning, and flag it as the
most important open question for follow-up work, ahead of further tuning our own method.

\section{Conclusion}

We implemented and evaluated a degree-aware curriculum hard-negative sampler for
LightGCL's BPR loss, verified on real Yelp2018 data that pure PPR hard-negative sampling
without protection measurably underperforms uniform sampling (supporting the
false-negative-bias concern that motivates this project), and showed that adding a
curriculum ramp recovers and exceeds the uniform baseline across overall metrics and
every degree bucket. Our full method's degree-gating component, however, did not show a
distinguishable benefit over a simpler fixed-ramp mixture in these experiments. We
consider the honest reporting of this partial result --- rather than only the parts that
worked --- the most useful outcome of the ablation, together with two numerical-stability
fixes to the reference LightGCL code (\S\ref{sec:stability}) that we believe are useful
independent of our sampler contribution.

\bibliographystyle{plain}
\begin{thebibliography}{9}
\bibitem{bpr} S. Rendle, C. Freudenthaler, Z. Gantner, and L. Schmidt-Thieme. BPR:
Bayesian personalized ranking from implicit feedback. In UAI, 2009.
\bibitem{lightgcn} X. He, K. Deng, X. Wang, Y. Li, Y. Zhang, and M. Wang. LightGCN:
Simplifying and powering graph convolution network for recommendation. In SIGIR, 2020.
\bibitem{pinsage} R. Ying, R. He, K. Chen, P. Eksombatchai, W. L. Hamilton, and J.
Leskovec. Graph convolutional neural networks for web-scale recommender systems. In KDD,
2018.
\bibitem{lightgcl} X. Cai, C. Huang, L. Xia, and X. Ren. LightGCL: Simple yet effective
graph contrastive learning for recommendation. In ICLR, 2023.
\bibitem{simgcl} J. Yu, H. Yin, X. Xia, T. Chen, L. Cui, and Q. V. H. Nguyen. Are graph
augmentations necessary? Simple graph contrastive learning for recommendation. In SIGIR,
2022.
\bibitem{curriculum} Y. Bengio, J. Louradour, R. Collobert, and J. Weston. Curriculum
learning. In ICML, 2009.
\end{thebibliography}

\appendix
\section{Note on dataset statistics discrepancy}
\label{app:dataset-note}
Our own project proposal's Table~1 reported Yelp2018 as 31{,}668 users / 38{,}048 items /
1{,}561{,}406 interactions (0.130\% density), sourced from the original LightGCN paper's
benchmark description. The data that actually ships with the official LightGCL
repository~\cite{lightgcl} (which the proposal itself specifies as the data source, "so
no extra preprocessing is required") has different statistics, reported in
Table~\ref{tab:dataset-stats}. We use the shipped data as-is, per the proposal's own
methodology, and report the measured statistics rather than the proposal draft's.

\section{Reproducibility}
Code, PPR precompute script, sampler implementation, screening/finals runner scripts, and
all raw per-run logs are in the project repository. Environment: PyTorch 2.11.0+cu128,
NVIDIA RTX 5090. Full commands and hyperparameter grids are in
\texttt{docs/design/sampler\_design.md} and \texttt{scripts/run\_screening.py} /
\texttt{scripts/run\_finals.py}.

\end{document}
